% Section "hist" (writer: hist). Source: research/tracks/history/notes-final.md (final numbering), with its
% referee.md and checks; foundations Prop 7.2 for lem:hist:chrono and cor:hist:H4. Credit to the author's notes
% (research/prior/hanni-sfi-notes.md) is in one footnote in the opening paragraph (REFRAME.md; log in
% research/paper-review/reframe/log-p1-hist-ex.md); no quotation of the notes.
% Proofs: sections/app-hist.tex.
% Renamings (NOTATION.md §4): tape measure \Unif; L_k -> \bar\ell_k; KH -> \Kpred; xi -> \MH; xi' -> \MHin;
% kappa -> \Kpad; tau -> \Ktail; (H1),(H2) -> (A1),(A2); private memory R -> \Gamma; computable random function
% pi -> nu; joint law mu -> \Prob; GF(2^m) -> GF(2^q); orders sigma, tau -> sigma, sigma'.
% Page cap (OUTLINE.md: "move secondary remarks ... to its appendix first"): the secondary items
% lem:hist:memory, prop:hist:nopathwise, rem:hist:thm1open, prop:hist:shannon, cor:hist:joint,
% lem:hist:chronolaw, prop:hist:variants and prop:hist:operator are stated in app-hist.tex; in the arXiv edit pass
% prop:hist:normchain, prop:hist:strict, conj:hist:algmi, prop:hist:lhg (with ass:hist:lhgeven, ass:hist:lhgodd),
% thm:hist:gf and thm:hist:expected moved there, and rem:hist:exploit was deleted (log-p1-hist-ex.md).
\section{History-aware prediction}\label{sec:hist}

SOI sees the past only as a set: its hypotheses answer each query without seeing the other inputs, and its loss $\LSFI(D_n)$ in the semimeasure convention ignores the order of the data (\cref{prop:setting:order}). We give two examples in which seeing the history helps without bound,\footnote{The examples of \cref{sec:hist:counter} and of \cref{thm:hist:randchain}, an informal dominance statement for simple reveal laws, the idea of proving it by simulating the reveal law inside a function machine (\cref{thm:hist:shared}) and the question behind \cref{conj:hist:algmi} first appeared in the author's working notes \citep{hanni2026notes}.} determine in which formulations a dominance statement for simple reveal laws holds (some cases remain open), and show in what sense the history-aware version of SOI, whose hypotheses see the list of previous inputs, is as good as SI (\cref{sec:hist:sfihist}). Lower bounds on $\LSFI$ are in the semimeasure convention unless stated; $c_\omega$ depends on the machines and on a Martin-L\"of random $\omega$. Proofs and some secondary statements are in \cref{app:hist}.

\subsection{History-aware predictors}\label{sec:hist:setting}

\begin{definition}[History-aware predictors]\src{history \S1.3, Lemma 1.0; foundations Prop 7.2}\label{def:hist:predictor}
A \emph{history-aware predictor} $P$ gives sub-probabilities $P(y\mid h,x)$; its loss is $L_P:=-\log P(y_{1:n}\Vert x_{1:n})$, $P(y_{1:n}\Vert x_{1:n}):=\prod_{i\le n}P(y_i\mid h_{i-1},x_i)$. A \emph{computable sampler} is a probabilistic machine that, given $(h,x)$, reads random bits and halts with a label or diverges (\emph{defective} if with positive probability); $P(y\mid h,x)$ is its probability of output $y$, lsc in the arguments. A sampler may keep memory between calls; it then predicts through its sequence law $P(y_{1:n}\Vert x_{1:n})$ and the conditionals this law induces (\cref{lem:hist:memory}). An \emph{input sampler} (reveal law) $\Delta$ outputs an input given $h$. A \emph{chronological semimeasure} is a map $\rho(y_{1:n}\Vert x_{1:n})\ge0$ with $\rho(\epsilon\Vert\epsilon)\le1$ and $\sum_y\rho(y_{1:n}y\Vert x_{1:n}x)\le\rho(y_{1:n}\Vert x_{1:n})$ for all $x$; $L_\rho:=-\log\rho(y_{1:n}\Vert x_{1:n})$.
\end{definition}

Chronological semimeasures go back to Hutter (see the paragraph after \cref{cor:hist:H4}). The dominance statements of \cref{sec:hist:thm1} concern computable samplers; one that halts with probability $1$ has computable probabilities ($1-\sum_{y'\ne y}P(y'\mid h,x)$ is upper semicomputable).

\begin{lemma}[SOI is a chronological semimeasure]\status{proved}\src{history Lemma 1.2; foundations Prop 7.2(b),(c)}\label{lem:hist:chrono}
$\rho_F(y_{1:n}\Vert x_{1:n}):=\Unif\{\omega:f_\omega(x_i)=y_i\ \forall i\le n\}$ is an lsc chronological semimeasure, equal to $\MF(D_n)$ on consistent sequences and $0$ otherwise. Hence the loss $\LH$ of \SFIhist{} (\cref{def:hist:sfihist}) satisfies $\LH\le\LSFI(D_n)+c$ in every order. The lsc chronological semimeasures that depend only on $D_n$ and vanish on inconsistent sequences are exactly the lsc weak function semimeasures.
\end{lemma}

\subsection{Counterexamples for order-free prediction}\label{sec:hist:counter}

In two examples a simple predictor that sees the history beats SOI without bound. In the first, a deterministic sequence of increasingly long random strings is revealed as a chain, each label being the previous input; in the second, every string of a given length but one is labelled $0$, and the remaining one, revealed last, is labelled $1$.

\begin{theorem}[The chain example]\status{proved}\src{history Thm 2.1}\label{thm:hist:chain}
Let $x_1,\dots,x_n$ ($n\ge2$) be distinct nonempty strings, $D_n:=\{(x_1,\epsilon)\}\cup\{(x_{i+1},x_i):i<n\}$, revealed as $(x_1,\epsilon),(x_2,x_1),\dots$.
(a) $K(x_1,\dots,x_{n-1}\mid x_n)-c\le\LSFI(D_n)\le K(x_1,\dots,x_{n-1})+c$ (semimeasure convention).
(b) If $1\le\ell_1<\ell_2<\cdots$ are computable, $x_i$ is the $i$-th block, of length $\ell_i$, of a Martin-L\"of random $\omega$, and $\bar\ell_k:=\sum_{i\le k}\ell_i$, then $\bar\ell_{n-1}-K(\ell_n)-c_\omega\le\LSFI(D_n)\le\bar\ell_{n-1}+K(n)+c$ (semimeasure convention); for $\ell_i=i$, $\LSFI(D_n)=n(n-1)/2\pm(2\log n+c_\omega)$.
(c) The predictor ``$\epsilon$ on the empty history; otherwise the unique input of the history that is not one of its labels'' has loss $0$; it is deterministic, set-based and of complexity $O(1)$.
\end{theorem}

\emph{Proof idea.} Iterating $f_\omega$ from $x_n$ recovers $x_{n-1},\dots,x_1$, so the chain events are disjoint (\cref{lem:hist:chain0}); apply the conditional coding theorem.

For the second example fix $s_\ell\in\{0,1\}^\ell$ for each $\ell$ and let $f_s(x):=\ind[x=s_{|x|}]$. Block $\ell$ reveals the other strings of length $\ell$ in lexicographic order, then $s_\ell$; $D^{(\ell)}$ is the data after blocks $1,\dots,\ell$.

\begin{theorem}[Skipped strings]\status{proved}\src{history Thm 2.2}\label{thm:hist:skipped}
(a) ``Predict $1$ at $x$ iff every other string of length $|x|$ is already an input'' has loss $0$; it is deterministic, set-based and of complexity $O(1)$. (b) $K(s_1,\dots,s_\ell\mid\ell)-c\le\LSFI(D^{(\ell)})\le\ell(\ell+1)/2+c$ (semimeasure convention). (c) If $s_1s_2\cdots$ are the consecutive blocks of a Martin-L\"of random $\omega$, then $\LSFI(D^{(\ell)})\ge\ell(\ell+1)/2-K(\ell)-c_\omega$; for each $\ell$, all but a $2^{-j}$ fraction of tuples $(s_1,\dots,s_\ell)$ give $\LSFI(D^{(\ell)})\ge\ell(\ell+1)/2-j-c$ (semimeasure convention).
\end{theorem}

The lower bound is the domain bound of \cref{thm:setting:coding}(a). In the length-lexicographic order \SFIhist{} pays as much as SOI (\cref{cor:hist:initial}), in the reveal order above, $O(1)$: the order carries the information.

\begin{proposition}[Every order-free predictor pays]\status{proved}\src{history Prop 2.3}\label{prop:hist:orderfree}
Let $\nu$ be a weak function semimeasure (\cref{def:setting:weak}), e.g.\ $\Mw$, $\Mprod$ or $W$. (a) For every $x$, $\sum_z\nu(D(x,z))\le1$ over all chains $D(x,z):=\{(x,z_1),(z_1,z_2),\dots,(z_k,\epsilon)\}$ ($k\ge0$) of distinct nonempty strings $z_i\ne x$. (b) If $\nu$ is lsc, the lower bounds of \cref{thm:hist:chain,thm:hist:skipped} hold for $-\log\nu$ with $K(\nu)$ subtracted. (c) For every $\nu$, with independent uniform blocks, $\Ex[-\log\nu(D_n)]\ge\bar\ell_{n-1}$ and $\Ex[-\log\nu(D^{(\ell)})]\ge\ell(\ell+1)/2$.
\end{proposition}

So every predictor that sees the past as a set \emph{and} is order-independent pays the cost (the predictors of \cref{thm:hist:chain}(c) and \cref{thm:hist:skipped}(a) are not order-independent). The normalised predictor $\MFbar$ is not order-free, but on both examples (with $\ell_i=i$ in the chain) it pays the same quadratic loss up to $O(n\log n)$, respectively $O(\ell\log\ell)$ bits (\cref{prop:hist:normchain}): a total \emph{broom} that names the past blocks and sends every other input to the last of them keeps its normaliser large (\cref{lem:hist:broom}).

\subsection{A dominance theorem for computable reveal laws}\label{sec:hist:thm1}

In these examples the reveal sequence is complex. If instead the inputs are drawn from a reveal law $\Delta$ that is simple to sample from given $f$, one may hope that every computable sampler $P$ that sees the history beats SOI in expected log loss by at most the complexity $C_\Delta$ of sampling from $\Delta$ plus a constant depending on $P$. How $C_\Delta$ is measured matters, since the randomness of the chain example can be moved into the reveal law (\cref{thm:hist:randchain}).

\begin{definition}[Three formulations]\src{history \S3}\label{def:hist:readings}
Let $y_i=f(x_i)$. \emph{Formulation A}: one reveal law $\Delta(x_i\mid h_{i-1})$ that does not see $f$ (private randomness and memory allowed), $f$ arbitrary. \emph{Formulation B} (the oracle form): $x_i\sim\Delta(\cdot\mid f,h_{i-1})$, with $C_\Delta$ the complexity of $\Delta$ relative to an oracle for $f$. \emph{Formulation C}: for each $f$ a reveal law $\Delta_f(x\mid h)$ that sees only the history, with $C_\Delta:=K(\Delta_f)$ as a standalone sampler; A is the case $\Delta_f=\Delta$.
\end{definition}

Formulation C is the natural one for the one-tape version of SOI: a function machine can simulate a reveal law that sees only the history, not one that sees $f$. The randomised chain has $C_\Delta=O(1)$ in Formulation B but is excluded in Formulation C (\cref{thm:hist:randchain}). Let $I(f\to x_{1:n}):=\sum_{i\le n}I(f;x_i\mid h_{i-1})$ (\emph{directed information}). A reveal law that does not see $f$ may keep private memory; it enters only through its \emph{reveal-sequence law} $\Delta(x_{1:n}\Vert y_{1:n-1})$, and $I(f\to x)=0$ (\cref{lem:hist:memory}). A natural proof simulates $\Delta$ and $P$ alternately inside a function machine until the query appears, outputs its label, and applies Jensen's inequality. With separate randomness per input, this direct simulation fails:

\begin{proposition}[A direct simulation fails]\status{computed; proved}\src{history Prop 3.1; check c1}\label{prop:hist:sketchfails}
Let $X=\{a,b\}$, $f(a)=f(b)=\beta$ for a uniform bit $\beta$, and let $\Delta$ show $a$, then $b$. The predictor with $y_1$ uniform and $y_2=y_1$ has loss $1$ bit; the direct simulator, run with separate randomness for each input, has loss $2$ bits.
\end{proposition}

A simulator with private randomness per input loses the cross-input correlations that $P$ uses, and the Jensen step averages over \emph{simulated} histories, whose labels come from $P$, not $f$. (The one-tape version itself pays at most $1+c$ bits here.) One simulation shared by all inputs works, at the price of the reveal sequence's log loss:

\begin{theorem}[Shared-randomness simulator]\status{proved}\src{history Thm 3.2}\label{thm:hist:shared}
For computable samplers $\Delta$ (inputs) and $P$ (labels), possibly defective and with private memory, and every consistent sequence, $\MF(D_n)\ge2^{-K(\Delta,P)-c}\Delta(x_{1:n}\Vert y_{1:n-1})P(y_{1:n}\Vert x_{1:n})$; so $\LSFI(D_n)\le L_P+L_\Delta+K(\Delta)+K(P)+c$ with $L_\Delta:=-\log\Delta(x_{1:n}\Vert y_{1:n-1})$. For distinct inputs the product may be summed over all orders of $D_n$.
\end{theorem}

\begin{theorem}[Bayes form]\status{proved}\src{history Thm 3.5; checks c2, r1}\label{thm:hist:bayes}
Let $f\sim\nu$, a computable random function (a machine law of a random total function, \cref{lem:setting:total}), and let the inputs see $f$ only through revealed labels: $x_i=\mathrm{sel}_i(\Gamma,h_{i-1})$, arbitrary maps, $\Gamma$ independent of $f$. Then the Bayes predictor $\nu(y\mid D_h,x)$ ($D_h$ the data of $h$; \cref{def:setting:data}) is the true conditional law, with loss $-\log\nu(D_n)$, and $\LSFI(D_n)\le-\log\nu(D_n)+K(\nu)+c$ on every path. So $\Ex\LSFI(D_n)\le\Ex L_P+K(\nu)+c$ for every predictor $P$, computable or not.
\end{theorem}

\begin{theorem}[Directed information]\status{proved; identity computed}\src{history Thm 3.6; checks c2, r1}\label{thm:hist:directed}
Let $f\sim\nu$, let the inputs have any joint law with $f$, $y_i=f(x_i)$, finite alphabets up to time $n$, and $L_\Prob:=-\sum_{i\le n}\log\Prob(y_i\mid h_{i-1},x_i)$. (a) For the law $\nu$ of $f$ (computable or not), $\Ex[-\log\nu(D_n)]-\Ex L_\Prob=I(f\to x_{1:n})-\Ex\KL(\Prob(f\in\cdot\mid h_n)\Vert\nu(f\in\cdot\mid f\supseteq D_n))\le I(f\to x_{1:n})$. (b) For computable $\nu$ and every $P$, $\Ex\LSFI(D_n)\le\Ex L_P+I(f\to x_{1:n})+K(\nu)+c$.
\end{theorem}

\begin{theorem}[The randomised chain; the oracle form fails]\status{proved; refutes Reading B}\src{history Thm 3.7; check r1}\label{thm:hist:randchain}
Let $x_i$ be independent and uniform on $\{0,1\}^i$, $f(x_{i+1}):=x_i$ and $f:=\epsilon$ elsewhere; reveal $x_1,x_2,\dots$. Then $f$ is a computable random function with $K(\nu)=O(1)$; an $O(1)$ program computes the reveal sequence from an oracle for $f$ ($x_i$ is the unique string of length $i\ge2$ with $f(x_i)\ne\epsilon$, and $x_1=f(x_2)$); the chain predictor has loss $0$; $I(f\to x_{1:n})=n(n+1)/2$; and $\Ex\LSFI(D_n)\ge n(n-1)/2$, $\Ex\LSFIbar\ge n(n-1)/2-n(2\log n+c)$.
\end{theorem}

This refutes only Formulation B: in Formulation C, $K(\Delta^{(n)}_f)\ge n(n+1)/2-j-c$ for all but a $2^{-j}$ fraction of chains, where $\Delta^{(n)}_f$ is the restriction of $\Delta_f$ to $n$ steps. One may ask for the statement with the algorithmic mutual information of $\Delta$ and $f$ in place of the complexity of $\Delta$; in Formulation B this also fails.

For arbitrary $f$ assume (A1) $Y$ is finite and (A2) a computable map $j\mapsto A_j$ to finite, listed subsets of $X$, of complexity $K(A)$, has all inputs up to step $n$ in $A_n$. Let $C_{\rm dir}(\Delta,n):=\sup_\nu I_\nu(f\to x_{1:n})$ over laws $\nu$ on $Y^{A_n}$ ($0$ if $\Delta$ does not see $f$).

\begin{theorem}[Minimax theorem]\status{(a) proved; (b) proved for reveal laws that do not see $f$, proof sketch for those that do; computed}\src{history Thm 3.8; checks c3, c8, r2}\label{thm:hist:minimax}
Fix $n$, assume (A1), (A2), and let the reveal law see the history (with private memory) and possibly $f|_{A_n}$ in case (a), and only the history in case (b). Assume that $P(y\mid h,x)$ and the reveal-sequence law have (a) rational values computable from their arguments, or (b) computable real values (this includes all computable samplers that halt with probability $1$; \cref{lem:hist:memory}). Then for every $f:A_n\to Y$, and on average over any law of $f$ independent of the reveal law's private randomness,
$\Ex_{x\mid f}\LSFI(D_n)\le\Ex_{x\mid f}L_P+C_{\rm dir}(\Delta,n)+K(P)+K(\Delta)+K(n)+K(A)+c$.
\end{theorem}

\emph{Proof idea.} Sion's theorem \citep{sion1958general} and \cref{thm:hist:directed}(a) give a law on $Y^{A_n}$ within $C_{\rm dir}+2$ bits of $P$ for every $f$, found by a search computable from $(P,\Delta,n,A)$.

\begin{remark}[Reveal laws that see $f$, case (b)]\status{proof sketch}\src{history Thm 3.8, Step 3, case (b)}\label{rem:hist:minimaxsketch}
For computable real values and reveal laws that may see $f|_{A_n}$, the argument of case (b) goes through if $C_{\rm dir}(\Delta,n)$ can be computed to within one bit from $(P,\Delta,n,A)$. \emph{Proof sketch.} $C_{\rm dir}$ is the maximum over a simplex of finitely many terms $u\log(u/v)$, with $u,v$ polynomial in the law of $f$ and in computable values of $\Delta$, and should be computable to within one bit by a modulus of continuity of entropy on a finite alphabet; with such a value the certified search in the proof of \cref{thm:hist:minimax} halts as before. We do not carry this out.
\end{remark}

\begin{corollary}[A dominance theorem for computable reveal laws]\status{proved under the stated hypotheses}\src{history Cor 3.9}\label{cor:hist:thm1}
(a) Under (A1), (A2), with $P$ and every $\Delta_f$ of computable real values (e.g.\ samplers halting with probability $1$, private memory allowed): $\Ex_{x\mid f}\LSFI(D_n)\le\Ex_{x\mid f}L_P+K(P)+K(\Delta_f)+K(n)+K(A)+c$ for every $f$; for random $f$, with $\Ex_fK(\Delta_f)$.
(b) If $x_i=\delta_f(h_{i-1})$ with $\delta_f$ computable: for every $f$, sampler $P$ (defective allowed) and $n$, $\LSFI(D_n)\le L_P+K(\delta_f)+K(P)+c$ on every path (\cref{cor:hist:deterministic}).
(c) For all samplers $P$ (defective allowed) and all reveal samplers $\Delta_f$ that halt with probability $1$, with no finiteness assumption, $\Ex_{x\mid f}\LSFI\le\Ex_{x\mid f}L_P+H(x_{1:n}\mid f)+K(\Delta_f)+K(P)+c$, where $H(x_{1:n}\mid f)$ is the entropy of the reveal sequence given $f$.
\end{corollary}

So in Formulation C the dominance statement holds pathwise for deterministic reveal policies (b); in expectation, for samplers that halt almost surely and finite alphabets up to the horizon, it holds with the additional term $K(n)+K(A)$ (a); and without finiteness assumptions it holds with the reveal entropy added (c). Defective samplers, infinite alphabets and the removal of $K(n)+K(A)$ are open (\cref{rem:hist:thm1open}); a direct simulation fails (\cref{prop:hist:sketchfails}). No pathwise version with $K(\Delta)$ alone holds (\cref{prop:hist:nopathwise}). The Shannon form of the mutual-information variant holds: if the reveal law is indexed by a random $\Theta$ and does not see $f$, then $I(f\to x)\le I(f;\Theta)$, with equality in the randomised chain (\cref{prop:hist:shannon}). The algorithmic form, with the algorithmic mutual information of $\Delta_f$ and $f$ in place of $K(\Delta_f)$, is open (\cref{conj:hist:algmi}).

\subsection{SOI with access to the previous inputs}\label{sec:hist:sfihist}

Between the versions of SOI considered so far and SI lies a natural intermediate, the history-aware version: the hypotheses see the list of previous inputs but not the previous labels, which they can recompute. We show in what sense it is as good as SI.

\begin{definition}[\SFIhist]\src{history \S4.1}\label{def:hist:sfihist}
Let $X_H$ be the set of nonempty finite input lists. A \emph{history-oracle function machine} is a function machine $T$ on $X_H$ (\cref{def:setting:machine}), queried at time $i$ at $x_{1:i}$; it never sees labels but can recompute its own. Its law $M^{\rm hist}_T(y_{1:n}\Vert x_{1:n}):=\Unif\{\omega:T^\omega(x_{1:i})=y_i\ \forall i\le n\}$ ($n\ge1$; $1$ at the root) is an lsc chronological semimeasure (\cref{lem:hist:chronolaw}). \emph{\SFIhist{}}, the history-aware version of SOI, is the one-tape version over $X_H$: its law is $\MH:=M^{\rm hist}_{T^H_U}$, $T^H_U$ the universal function machine for $X_H$ (\cref{thm:setting:universal}), and its loss is $\LH:=-\log\MH(y_{1:n}\Vert x_{1:n})$. We \emph{define} the universal chronological semimeasure as $\MH$.
\end{definition}

\begin{theorem}[Representation]\status{proved; construction computed}\src{history Thm 4.2; checks c4, r5}\label{thm:hist:representation}
For every lsc chronological semimeasure $\rho$ there is a history-oracle function machine $T_\rho$ with $M^{\rm hist}_{T_\rho}=\rho$ at every non-root node and $K(T_\rho)\le K(\rho)+c$.
\end{theorem}

The proof allocates dyadic intervals as \citet{zvonkin1970complexity} do for sequences; for sequences, \citet{wood2013non} show that machine-based and semimeasure-based universal priors form the same class.

\begin{corollary}[Universality of \SFIhist]\status{proved; identification known (unverified)}\src{history Cor 4.3; foundations Prop 7.2(a),(d)}\label{cor:hist:H4}
(a) $\MH\ge2^{-K(\rho)-c}\rho$ for every lsc chronological semimeasure $\rho$; so $\MH$ is a universal element of this class and agrees with every other one up to constant factors.
(b) For every lsc history-aware predictor $P$ (e.g.\ every computable sampler, defective or not) and \emph{every} sequence, consistent or not: $\LH\le L_P+K(P)+c$.
(c) $\LH\le\LSFI(D_n)+c$ on every consistent sequence, in every order.
\end{corollary}

In this sense \SFIhist{} is as good as SI. A trimming construction as in \cref{thm:setting:weakuniversal} gives the same law up to constants. Chronological semimeasures and a universal mixture over them were introduced by Hutter for the AIXI model \citep{hutter2000theory,hutter2005universal}, with actions and percepts; here the inputs play the actions and the labels the percepts. We define $\MH$ directly and do not use any identification with Hutter's construction. \citet{hutter2007aixi} lists supervised learning among the problem classes of the AIXI model. Oracles that see the previous pairs, the time $i$, or the set of previous inputs define \SFIpairs, \SFIcount{} (SOI on time-stamped inputs) and \SFIset, with laws $\Mpairs,\Mcount,\Mset$ and losses $L_{\rm pairs},L_{\rm count},L_{\rm set}$; on distinct inputs SOI $\preceq$ \SFIcount{} $\preceq$ \SFIset{} $\preceq$ \SFIhist{} up to constant factors, and $\Mpairs$ agrees with $\MH$ up to constant factors (\cref{prop:hist:variants}). Each inclusion is strict, with gaps at least linear in $n$ on suitable orders and labels and also along single infinite sequences (semimeasure convention; \cref{prop:hist:strict,prop:hist:strictinf}); so the oracle must give the list of previous inputs, not their set (also \cref{prop:ex:cycset}).

Solomonoff's per-query solution $\Mpq$ (\cref{rem:versions:operator}) ignores the history: \SFIhist{} dominates it, and on distinct inputs it pays the lower bounds of \cref{thm:hist:chain,thm:hist:skipped} (\cref{prop:hist:operator}).

For deterministic programs on $X_H$, the Kolmogorov and mixture versions differ unboundedly, as for SOI (\cref{thm:versions:nocoding}): against fair-coin labels the program mixture pays at most $\Ktail(n)+c$ extra bits, which grows more slowly than every computable unbounded function, while for the worst labelling, along every input sequence, the single shortest program pays more than $\log n+\log\log n$ extra bits infinitely often (\cref{prop:hist:programs}).

\subsection{As good as Solomonoff induction?}\label{sec:hist:si}

Let $M$ be the universal monotone semimeasure (\cref{prop:setting:si}), $\code(h_n):=\bar x_1\bar y_1\cdots\bar x_n\bar y_n$ with $\bar s:=1^{|s|}0s$, and $\LMjoint:=-\log M(\code(h_n))$; SOI's loss is at most $\LMjoint+c$ (\cref{cor:hist:joint}). Let $\LMx$ be the sum over $i$ of $-\log$ of $M$'s conditional probability of $\bar x_i$ given $\code(h_{i-1})$ (the loss on the inputs), and $\LMy:=\LMjoint-\LMx$, Solomonoff induction's unnormalised conditional label loss. Let $\MHin$ be the law of the universal function machine on finite label lists (the empty list included) that answers with the next input; by \cref{thm:hist:representation} with the roles exchanged (same proof), $\MHin\ge2^{-K(\rho')-c}\rho'$ for every lsc chronological semimeasure $\rho'(x_{1:n}\Vert y_{1:n-1})$ for inputs given previous labels. Let $L^x_{\rm in}:=-\log\MHin(x_{1:n}\Vert y_{1:n-1})$.

\begin{theorem}[Solomonoff induction on the joint sequence]\status{proved}\src{history Thm 4.6, Lemma B}\label{thm:hist:si}
On every sequence: (a) $\LH\le\LMjoint+c$ and $L^x_{\rm in}\le\LMjoint+c$; (b) $\LMjoint\le\LH+L^x_{\rm in}+c$; (c) $\LH-\LMy\le\LMx+c$ and $\LMy-\LH\le L^x_{\rm in}-\LMx+c$; (d) if $x_i=\delta(h_{i-1})$ for a partial computable $\delta$, then $\LMjoint\le\LH+K(\delta)+c$. This holds for every prefix-free block code with nonempty codewords.
\end{theorem}

\begin{proposition}[Input-loss identity]\status{proved; toy computed}\src{history Prop 4.7; check s3}\label{prop:hist:inputloss}
For a partial computable $\delta$ and every sequence with $x_i=\delta(h_{i-1})$ for all $i$: $\LMx-K(\delta)-c\le\LH-\LMy\le\LMx+c$.
\end{proposition}

So for deterministic input policies, ``$\LH\le\LMy+O(1)$'' holds iff $M$ predicts the inputs with bounded loss. Under two assumptions on unnormalised $M$ that are suggested by work of Lattimore, Hutter and Gavane \citep{lattimore2011universal} and are not proved here (\cref{ass:hist:lhgeven,ass:hist:lhgodd}), neither $\LMy\le\LH+O(1)$ nor $\LH\le\LMy+O(1)$ holds on all sequences of binary inputs and labels under the identity code (\cref{prop:hist:lhg}); this does not contradict \cref{cor:hist:H4}, as $M$'s conditionals need not be lsc. For normalised $M$ both directions are open.

\begin{remark}[Summary: hypotheses that see the previous inputs]\status{remark (cited items)}\src{history \S7 ``H4''}\label{rem:hist:H4verdict}
With the list of previous inputs SOI becomes the universal chronological semimeasure (\cref{cor:hist:H4}); the set is not enough (\cref{prop:hist:strict}(iii)). \SFIhist{} is as good as SI against Solomonoff induction's joint loss (\cref{thm:hist:si}(a)); against its conditional label loss the difference is $M$'s loss on the inputs (\cref{prop:hist:inputloss}), and for unnormalised $M$ it fails both ways, conditionally on Assumptions LHG-even and LHG-odd (\cref{prop:hist:lhg}); the normalised case is open.
\end{remark}

\subsection{What SOI can and cannot exploit}\label{sec:hist:exploit}

For $D$ with distinct inputs let $D_\sigma$ be $D$ in order $\sigma$, $x_\sigma$ its input list, and $\Kpred(x_{1:n}):=\min\{|p|:V(p,x_i)=\langle x_1,\dots,x_{i-1}\rangle\ \forall i\le n\}$, the complexity of recovering each input's predecessors from the input.

\begin{theorem}[Sandwich]\status{proved}\src{history Thm 5.2 (with Prop 5.1)}\label{thm:hist:sandwich}
For every data set $D$: $\max_\sigma\LH(D_\sigma)-c\le\LSFI(D)\le\min_{\sigma'}[\LH(D_{\sigma'})+\Kpred(x_{\sigma'})]+c$.
\end{theorem}

\begin{corollary}[Initial segments]\status{proved}\src{history Cor 5.3}\label{cor:hist:initial}
If $\dom D=\{e_1,\dots,e_k\}$ for a computable injective sequence $e$, then $|\LSFI(D)-\LH(D\text{ in $e$-order})|\le K(e)+c$, and the $e$-order is the worst for \SFIhist{} up to $K(e)+c$.
\end{corollary}

This covers the skipped strings and the full cyclic permutation in the lexicographic order of $\{0,1\}^d$ ($K(e)\le K(d)+c$); in its reveal order, \SFIhist{} pays $O(1)$ (\cref{prop:ex:cychist}). The $\Kpred$ term can be needed beyond order information. For $a\ne b$ and $r$ in $\{0,1\}^q$, read as elements of $\mathrm{GF}(2^q)$, and $D_{abr}:=\{(a,r),(b,r+ab)\}$, \SFIhist{} pays at most $q+c$ in both orders, as it predicts $y_2=y_1+x_1x_2$, while SOI, whose hypotheses never see the first input, pays at least $3q/2-2$ bits in expectation over uniform triples with $a\ne b$ and at least $3q/2-j-c$ bits on all but a $2^{1-j}$ fraction of them (\cref{thm:hist:gf}).

In expectation, for $f$ drawn from a computable law $\nu$, inputs drawn by any kernel that may see $f$ and finite alphabets, the gap $\Ex\LSFI(D_n)-\Ex\LH$ lies between $-c$ and $I(f\to x_{1:n})+K(\nu)+c$, and is at most $K(\nu)+c$ if the input choice does not see $f$; the upper bound is attained up to the last input's entropy in the randomised chain, but not in general (\cref{thm:hist:expected}).

On individual sequences, $\LSFI(D)\le\LH(D_\sigma)+I(x_\sigma:D)+c$, with $I(a:b):=K(b)-K(b\mid a)$; the natural strengthening with $I(x_\sigma:y_\sigma)$ in place of $I(x_\sigma:D)$ is false, by a margin linear in $n$ (\cref{prop:hist:info}). Informally, shared randomness lets a label of the one-tape version correlate with a \emph{guess} about the other inputs (as the brooms and the optimal law for \cref{thm:hist:gf} do), but not depend on which inputs are present.

\begin{question}[history-aware prediction]\label{q:hist:open}
(i) In Formulation C (\cref{rem:hist:thm1open}): can $K(n)+K(A)$ be dropped, and do defective samplers, infinite alphabets and, in case (b) of \cref{thm:hist:minimax}, reveal laws that see $f$ admit the same bound? (ii) Do Assumptions LHG-even and LHG-odd hold, and how does $\LH$ compare with the label loss of normalised $M$? (iii) Do the lower bounds of \cref{prop:hist:strict,prop:hist:strictinf,thm:hist:gf} hold for normalised SOI? (iv) Is $\LMjoint\ge\LH+L^x_{\rm in}-O(1)$ (the converse is \cref{thm:hist:si}(b))? (v) How does $\max_GN(G)$ in \cref{lem:hist:charsum} grow with $q$? (vi) Does \cref{conj:hist:algmi} hold? (vii) Which individual-sequence quantity measures $\LSFI(D)-\LH(D_\sigma)$ up to $O(\log n)$? (viii) What is the exact growth rate of the excess in \cref{prop:hist:programs}(d)?
\end{question}
